\subsection{Sequential Containers}

\subsubsection{Element Access}

\begin{tabularx}{\linewidth}{l|c|c|c|c}
& \textbf{array} & \textbf{vector} & \textbf{deque} & \textbf{list} \\
\hline
\texttt{[]} & \checkmark & \checkmark & \checkmark & $\times$ \\
\texttt{at()} & \checkmark & \checkmark & \checkmark & $\times$ \\
\texttt{front()} & \checkmark & \checkmark & \checkmark & \checkmark \\
\texttt{back()} & \checkmark & \checkmark & \checkmark & \checkmark \\
\texttt{data()} & \checkmark & \checkmark & $\times$ & $\times$ \\
\end{tabularx}

\subsubsection{Array}

\textbf{Definition:}\\
\texttt{array<T, N> a;}\\

Fixed-size container. Size must be known at compile time.

\textbf{Modifiers}

\begin{tabularx}{\linewidth}{lX}
\texttt{a.fill(value)} & Fill array with value.\\
\end{tabularx}

\subsubsection{Vector}

\textbf{Definition:}\\
\texttt{vector<T> v;}

Dynamic array with automatic resizing. Supports all array methods plus:

\textbf{Additional Modifiers}

\begin{tabularx}{\linewidth}{lX}
\texttt{v.push\_back(x)} & Add element to the end.\\
\texttt{v.pop\_back()} & Remove last element.\\
\texttt{v.emplace\_back(args)} & Construct element in place at the end.\\
\texttt{v.insert(pos, x)} & Insert element at iterator position.\\
\texttt{v.emplace(pos, args)} & Construct element in place at position.\\
\texttt{v.erase(pos)} & Remove element at iterator position.\\
\end{tabularx}

\textbf{Capacity Management}

\begin{tabularx}{\linewidth}{lX}
\texttt{v.capacity()} & Current allocated storage.\\
\texttt{v.reserve(n)} & Reserve storage for n elements.\\
\texttt{v.shrink\_to\_fit()} & Reduce capacity to fit size.\\
\end{tabularx}

\vspace*{8pt}
\texttt{tie} can be used to unpack a vector.\\
\textit{Example:} \texttt{tie(a, b, c) = v;}


\vfill \null
\columnbreak


\subsubsection{Deque}

\textbf{Definition:}\\
\texttt{deque<T> d;}

Double-ended queue. Supports all array element access (except \texttt{data()}) and all vector modifiers (except capacity management), plus:

\textbf{Additional Front Operations}

\begin{tabularx}{\linewidth}{lX}
\texttt{d.push\_front(x)} & Add element to the front.\\
\texttt{d.pop\_front()} & Remove first element.\\
\texttt{d.emplace\_front(args)} & Construct element in place at the front.\\
\end{tabularx}
